\subsection{Радиус кривизны, угловое ускорение}
%%%%%%%%%%%%%%%%%%%%%%%%%%%%%%%%%%%%%%%%%%%%%%%%%%%%
%%%%%%%%%%%%%%%%%%%%%%%%%%%%%%%%%%%%%%%%%%%%%%%%%%%%
\z{Тело бросают со скоростью $v_0$ под углом $\alpha$ к горизонту. Найдите зависимость радиуса кривизны траектории от времени $R_\text{кр}(t)$.}
%
{$R_\text{кр}(t)=\frac{\left((v_0\cos{\alpha})^2+(v_0\sin\alpha-gt)^2\right)^\frac{3}{2}}{gv_0\cos{\alpha}}$}

%%%%%%%%%%%%%%%%%%%%%%%%%%%%%%%%%%%%%%%%%%%%%%%%%%%%

 \z[{\begin{wrapfigure}[3]{r}{0.3\textwidth}%
		\vspace{-15pt}%
		\begin{flushright}%
 			\includegraphics[width=0.3\textwidth]{3.2.2}%
		\end{flushright}%
	\end{wrapfigure}%
	}]%
	{Край гладкого горизонтального стола скруглен по окружности радиуса $r$. С какой наименьшей скоростью нужно пустить по столу малое тело, чтобы оно, достигнув скругления, сразу полетело по параболе? }
%
{$v=\sqrt{gr}$}
%%%%%%%%%%%%%%%%%%%%%%%%%%%%%%%%%%%%%%%%%%%%%%%%%%%%
\z{Сферический резервуар, стоящий на земле, имеет радиус $R$. При какой наименьшей скорости брошенный с земли камень может перелететь через резервуар, лишь коснувшись вершины?}
%
{$v=\sqrt{5gr}$}

%%%%%%%%%%%%%%%%%%%%%%%%%%%%%%%%%%%%%%%%%%%%%%%%%%%%	
\z{В момент времени, когда скорость частицы равна $10^6\,  \text{м/с}$, её ускорение составляет $10^4\,  \text{м}/\text{с}^2$ и направлено под углом $30\degree$ к скорости. На сколько увеличится скорость за $10^{-2}\, \text{с}$? На какой угол изменится направление скорости? Какова в этот момент угловая скорость вращения вектора скорости?}
%
{\puan$\frac{\sqrt{3}}{2}\cdot 10^2 \,\text{м/с}$;~\puan$5\cdot10^{-5}\,\text{рад}$;~\puan$\omega = 5\cdot 10^{-3}\,\text{с}^{-1}$}

%%%%%%%%%%%%%%%%%%%%%%%%%%%%%%%%%%%%%%%%%%%%%%%%%%%%
\z{Кусочек сахара брошен со скоростью $v_0$ под углом $\alpha$ к горизонту. Вслед за сахаром точно по той же траектории летит пчела с постоянной по модулю скоростью $v_0$. Найдите ускорение пчелы в верхней точке.}
%
{$a=\frac{g}{\cos^2{\alpha}}$}

